% GENERATED FILE. Edit canonical lessons, then run npm run generate:books.
% canonical-source-hash: fnv1a64:59b134b1e812c298
% canonical-lessons: LA-C206-anima, LA-C206-memoria, LA-C206-sapientia, LA-C206-ingenium, LA-C206-consuetudo

\chapter{Breath, Memory, Wisdom, Talent, Habit}
\label{ch:la-breath-memory-wisdom-talent-habit}
\hlchaptermodality{206}

\begin{chapteropening}
\textbf{By the end of this chapter:} \emph{I can say breath, memory, wisdom, talent and a habit.}

Complete the last lesson of chapter 206: I can say breath, memory, wisdom, talent and a habit.
\end{chapteropening}

\section[anima, animae]{anima, animae — breath, soul}

\label{lesson:LA-C206-anima}



\begin{quote}
Before the new one: say the Latin for a drug, a medicine, then the Latin for an ointment.
\end{quote}

\subsection*{You'll want to know: anima, animae}
\textbf{anima, animae} — \textquotedblleft{}breath, soul\textquotedblright{}.

\begin{cousinweb}
\textbf{anima} is the partner of \emph{animus}, \textquotedblleft{}mind, spirit\textquotedblright{} (EtymDB 2.1), and both are usually traced to Proto-Indo-European \emph{*h2enh1-}, \textquotedblleft{}to breathe\textquotedblright{} (EtymDB 2.1): the soul as breath.
\end{cousinweb}

\begin{grammarlens}[title={What you've built}]
One more word for this chapter.
\end{grammarlens}

\subsection*{Your turn}
\begin{itemize}
\raggedright
  \item \emph{Say it:} \emph{anima, animae}
  \item \emph{Say it:} \emph{anima, animae}, once more
  \item \emph{Say it:} say \emph{anima, animae}
  \item \emph{Recall:} say the Latin for a drug, a medicine, then the Latin for an ointment, then say \emph{anima, animae} again
\end{itemize}

\begin{tcolorbox}[breakable,colback=teal!4,colframe=teal!35!black,title={Before you move on}]
What does anima, animae mean? (\textquotedblleft{}Breath, soul\textquotedblright{}.) Say it once more.
\end{tcolorbox}
\section[memoria, memoriae]{memoria, memoriae — memory}

\label{lesson:LA-C206-memoria}



\begin{quote}
Before the new one: say the Latin for an ointment, then the Latin for breath.
\end{quote}

\subsection*{You'll want to know: memoria, memoriae}
\textbf{memoria, memoriae} — \textquotedblleft{}memory\textquotedblright{}.

\begin{cousinweb}
\textbf{memoria} is built on \emph{memor}, \textquotedblleft{}mindful\textquotedblright{} (Lewis \& Short), usually traced to Proto-Indo-European \emph{*(s)mer-}, \textquotedblleft{}to remember, to be mindful\textquotedblright{} (EtymDB 2.1).
\end{cousinweb}

\begin{grammarlens}[title={What you've built}]
One more word for this chapter.
\end{grammarlens}

\subsection*{Your turn}
\begin{itemize}
\raggedright
  \item \emph{Say it:} \emph{memoria, memoriae}
  \item \emph{Say it:} \emph{memoria, memoriae}, once more
  \item \emph{Say it:} say \emph{memoria, memoriae}
  \item \emph{Recall:} say the Latin for an ointment, then the Latin for breath, then say \emph{memoria, memoriae} again
\end{itemize}

\begin{tcolorbox}[breakable,colback=teal!4,colframe=teal!35!black,title={Before you move on}]
What does memoria, memoriae mean? (\textquotedblleft{}Memory\textquotedblright{}.) Say it once more.
\end{tcolorbox}
\section[sapientia, sapientiae]{sapientia, sapientiae — wisdom}

\label{lesson:LA-C206-sapientia}



\begin{quote}
Before the new one: say the Latin for breath, then the Latin for memory.
\end{quote}

\subsection*{You'll want to know: sapientia, sapientiae}
\textbf{sapientia, sapientiae} — \textquotedblleft{}wisdom\textquotedblright{}.

\begin{cousinweb}
\textbf{sapientia} is built on \emph{sapiēns}, \textquotedblleft{}wise\textquotedblright{} (Lewis \& Short), from \emph{sapiō}, \textquotedblleft{}to taste, to be wise\textquotedblright{}, with \emph{-ia}, an ending for a quality.
\end{cousinweb}

\begin{grammarlens}[title={What you've built}]
One more word for this chapter.
\end{grammarlens}

\subsection*{Your turn}
\begin{itemize}
\raggedright
  \item \emph{Say it:} \emph{sapientia, sapientiae}
  \item \emph{Say it:} \emph{sapientia, sapientiae}, once more
  \item \emph{Say it:} say \emph{sapientia, sapientiae}
  \item \emph{Recall:} say the Latin for breath, then the Latin for memory, then say \emph{sapientia, sapientiae} again
\end{itemize}

\begin{tcolorbox}[breakable,colback=teal!4,colframe=teal!35!black,title={Before you move on}]
What does sapientia, sapientiae mean? (\textquotedblleft{}Wisdom\textquotedblright{}.) Say it once more.
\end{tcolorbox}
\section[ingenium, ingeniī]{ingenium, ingeniī — talent, character}

\label{lesson:LA-C206-ingenium}



\begin{quote}
Before the new one: say the Latin for memory, then the Latin for wisdom.
\end{quote}

\subsection*{You'll want to know: ingenium, ingeniī}
\textbf{ingenium, ingeniī} — \textquotedblleft{}talent, character\textquotedblright{}.

\begin{cousinweb}
\textbf{ingenium} is \emph{in}, \textquotedblleft{}in\textquotedblright{}, with the root of \emph{gignō}, \textquotedblleft{}to beget, to give birth to\textquotedblright{} (Lewis \& Short): what is born in a person, one's nature. English \emph{engine} and \emph{ingenious} come from it.
\end{cousinweb}

\begin{grammarlens}[title={What you've built}]
One more word for this chapter.
\end{grammarlens}

\subsection*{Your turn}
\begin{itemize}
\raggedright
  \item \emph{Say it:} \emph{ingenium, ingeniī}
  \item \emph{Say it:} \emph{ingenium, ingeniī}, once more
  \item \emph{Say it:} say \emph{ingenium, ingeniī}
  \item \emph{Recall:} say the Latin for memory, then the Latin for wisdom, then say \emph{ingenium, ingeniī} again
\end{itemize}

\begin{tcolorbox}[breakable,colback=teal!4,colframe=teal!35!black,title={Before you move on}]
What does ingenium, ingeniī mean? (\textquotedblleft{}Talent, character\textquotedblright{}.) Say it once more.
\end{tcolorbox}
\section[cōnsuētūdō, cōnsuētūdinis]{cōnsuētūdō, cōnsuētūdinis — a habit, a custom}

\label{lesson:LA-C206-consuetudo}



\begin{quote}
Before the new one: say the Latin for wisdom, then the Latin for talent.
\end{quote}

\subsection*{You'll want to know: cōnsuētūdō, cōnsuētūdinis}
\textbf{cōnsuētūdō, cōnsuētūdinis} — \textquotedblleft{}a habit, a custom\textquotedblright{}.

\begin{cousinweb}
\textbf{cōnsuētūdō} is built on \emph{cōnsuēscō}, \textquotedblleft{}to get used to\textquotedblright{} (Lewis \& Short), with \emph{-tūdō}, an ending for a state. English \emph{custom} comes from it through French.
\end{cousinweb}

\begin{grammarlens}[title={What you've built}]
One more word for this chapter.
\end{grammarlens}

\subsection*{Your turn}
\begin{itemize}
\raggedright
  \item \emph{Say it:} \emph{cōnsuētūdō, cōnsuētūdinis}
  \item \emph{Say it:} \emph{cōnsuētūdō, cōnsuētūdinis}, once more
  \item \emph{Say it:} say \emph{cōnsuētūdō, cōnsuētūdinis}
  \item \emph{Recall:} say the Latin for wisdom, then the Latin for talent, then say \emph{cōnsuētūdō, cōnsuētūdinis} again
\end{itemize}

\begin{tcolorbox}[breakable,colback=teal!4,colframe=teal!35!black,title={Before you move on}]
What does cōnsuētūdō, cōnsuētūdinis mean? (\textquotedblleft{}A habit, a custom\textquotedblright{}.) Say it once more.
\end{tcolorbox}
